% AgentCanary攻击类型与本研究风险类型对应关系表
\begin{table}[htbp]
\centering
\caption{AgentCanary攻击类型与本研究$R_0$--$R_6$风险类型对应关系}
\label{tab:agentcanary_risk_mapping}
\small
\begin{tabularx}{\textwidth}{clXl}
\toprule
\textbf{风险类型} & \textbf{定义} & \raggedright\arraybackslash\textbf{对应AgentCanary攻击类型} & \textbf{对应风险入口} \\
\midrule
$R_0$ & 正常或已阻断 & no-attack（无攻击基线） & 全部入口 \\
$R_1$ & 过度授权/越权操作 & permission\_escalation, defense\_disable & DPI、IPI \\
$R_2$ & 敏感数据外传 & privacy\_leak, phishing & DPI、IPI \\
$R_3$ & 不可逆操作缺少确认 & file\_deletion, resource\_exhaustion & IPI、IF \\
$R_4$ & 工具参数来源污染 & env\_tampering, integrity\_tamper & IPI、SP \\
$R_5$ & 跨Agent委托/权限失控 & unauthorized\_access, privilege\_inheritance & SP \\
$R_6$ & 共享记忆/持久状态污染 & illegal\_memory, persona\_corruption & MC \\
\bottomrule
\end{tabularx}
\end{table}
